\chapter{Оптимизация образов}

Маленький образ быстрее скачивается, быстрее запускается, занимает
меньше места в реестре и оставляет меньше поверхности атаки. При этом
оптимизация не должна превращаться в самоцель: важнее
воспроизводимость и понятность. В главе~--- что именно занимает место,
как выбрать базу и как измерить результат.

\section{Почему размер важен}

\begin{itemize}
  \item \textbf{Скорость доставки.} Образ в 1~ГБ скачивается
        в десять раз дольше, чем в 100~МБ.
  \item \textbf{Холодный старт.} На новом узле контейнер не запустится,
        пока не распакуются слои.
  \item \textbf{Стоимость хранения.} Реестры и диски узлов не бесплатны.
  \item \textbf{Поверхность атаки.} Каждый лишний пакет~--- потенциальная
        уязвимость.
  \item \textbf{Скорость CI.} Меньше данных~--- быстрее \cmd{push}
        и \cmd{pull} в конвейере.
\end{itemize}

\begin{notebox}
Размер не единственная метрика. Образ на 300~МБ с понятным
Dockerfile лучше образа на 150~МБ, собранного хрупким скриптом.
Сначала правильность и воспроизводимость, потом байты.
\end{notebox}

\section{Что занимает место}

Посмотрите на слои и их размеры.

\begin{code}[language=bash]
$ docker image ls
$ docker history --human --no-trunc api:1.0
$ docker system df -v
\end{code}

Типичные источники лишнего веса:

\begin{itemize}
  \item полный базовый образ с компиляторами и документацией;
  \item кэши пакетных менеджеров (\cmd{\textasciitilde/.cache/pip},
        \cmd{/root/.npm}, \cmd{/var/lib/apt/lists});
  \item инструменты сборки, оставшиеся в финальном слое;
  \item исходники и тесты, скопированные \cmd{COPY . .};
  \item отладочные символы в бинарниках;
  \item несколько пакетных менеджеров в одном образе.
\end{itemize}

Для визуального анализа удобен \cmd{dive}: он показывает слои,
их размер и \enquote{мусор} внутри.

\begin{code}[language=bash]
$ dive api:1.0
\end{code}

\section{Выбор базового образа}

\begin{center}
\small
\begin{tabularx}{\textwidth}{@{}lXX@{}}
\toprule
\textbf{База} & \textbf{Размер} & \textbf{Особенности} \\
\midrule
\cmd{scratch} & 0 & пустой; только статические бинарники, нет
оболочки и \file{/etc/passwd} \\
\cmd{distroless} & 2--20 МБ & минимальные библиотеки, нет shell;
удобно для Go, Java, Python с осторожностью \\
\cmd{alpine} & 5--10 МБ & musl libc, \cmd{apk}; экономия, но возможны
сюрпризы с бинарными wheel'ами \\
\cmd{-slim} & 50--90 МБ & Debian без лишнего; привычный glibc
и \cmd{apt} \\
полный образ & 150--900 МБ & компиляторы, заголовки, утилиты; годится
как builder, но не как runtime \\
\bottomrule
\end{tabularx}
\end{center}

\begin{warnbox}
Alpine экономит десятки мегабайт, но использует musl вместо glibc.
Некоторые Python-пакеты и JVM-библиотеки требуют сборки или ведут
себя иначе, а DNS и locale отличаются от Debian. Проверяйте
приложение на выбранной базе, а не только на \cmd{ubuntu}.
\end{warnbox}

\section{Порядок слоёв и очистка}

\begin{code}[language=dockerfile]
# Плохо: кэш apt остаётся в слое
RUN apt-get update
RUN apt-get install -y curl
RUN rm -rf /var/lib/apt/lists/*

# Хорошо: всё в одном слое
RUN apt-get update \
 && apt-get install -y --no-install-recommends curl \
 && rm -rf /var/lib/apt/lists/*
\end{code}

Удаление файла в следующем слое не уменьшает образ: файл остаётся
в нижнем слое, а сверху появляется whiteout. Чистить нужно в том же
\cmd{RUN}, где создан мусор.

Для Python и Node.js:

\begin{code}[language=dockerfile]
RUN pip install --no-cache-dir -r requirements.txt
RUN npm ci --omit=dev && npm cache clean --force
\end{code}

\section{Кэш сборки}

Кэш не влияет на размер финального образа, но сильно влияет на время
сборки.

\begin{code}[language=dockerfile]
COPY requirements.txt .
RUN --mount=type=cache,target=/root/.cache/pip \
    pip install -r requirements.txt
COPY . .
\end{code}

Кэш-монтирование сохраняет скачанные пакеты между сборками, не попадая
в образ. Для Node.js аналогично используется \cmd{/root/.npm}.

\begin{tipbox}
В CI включайте \flag{--cache-from type=registry} или
\flag{--cache-to type=gha}. Это даёт ускорение в разы по сравнению
с чистой сборкой на каждом коммите.
\end{tipbox}

\section{Многоэтапная сборка как основной приём}

Идея проста: тяжёлые инструменты нужны только для сборки, а в runtime
попадает результат.

\begin{code}[language=dockerfile]
FROM node:22-alpine AS deps
WORKDIR /app
COPY package*.json ./
RUN npm ci --omit=dev

FROM node:22-alpine
WORKDIR /app
COPY --from=deps /app/node_modules ./node_modules
COPY --chown=node:node . .
USER node
ENV NODE_ENV=production
CMD ["node", "server.js"]
\end{code}

Для Go финальный слой может быть \cmd{scratch}:

\begin{code}[language=dockerfile]
FROM golang:1.25-alpine AS build
WORKDIR /src
COPY go.mod go.sum ./
RUN go mod download
COPY . .
RUN CGO_ENABLED=0 go build -trimpath -ldflags="-s -w" -o /app ./cmd/app

FROM scratch
COPY --from=build /app /app
USER 10001
ENTRYPOINT ["/app"]
\end{code}

Флаг \flag{-ldflags="-s -w"} убирает отладочные символы,
\flag{-trimpath}~--- абсолютные пути сборки. Вместе они дают
заметную экономию и немного улучшают воспроизводимость.

\section{Пример: от 1 ГБ к 150 МБ}

\begin{code}[language=dockerfile]
# Было
FROM python:3.13
COPY . .
RUN pip install -r requirements.txt
CMD ["python", "app.py"]
\end{code}

Проблемы: полный образ с компиляторами, весь контекст (включая тесты
и \file{.git}) в образе, кэш pip, запуск от root, нет закреплённых
зависимостей.

\begin{code}[language=dockerfile]
# Стало
FROM python:3.13-slim

ENV PYTHONDONTWRITEBYTECODE=1 \
    PYTHONUNBUFFERED=1

WORKDIR /app
COPY requirements.txt .
RUN pip install --no-cache-dir -r requirements.txt
COPY --chown=10001:10001 app.py .
USER 10001
CMD ["python", "app.py"]
\end{code}

Плюс \file{.dockerignore}:

\begin{code}
.git
tests
__pycache__
*.pyc
.env
.venv
\end{code}

Результат: образ уменьшается примерно с 1~ГБ до 150~МБ, а слой
с зависимостями переиспользуется при изменении \file{app.py}.

\begin{code}[language=bash]
$ docker image ls api
REPOSITORY   TAG      IMAGE ID       SIZE
api          before   a11111111111   1.02GB
api          after    b22222222222   152MB
\end{code}

\section{Ускорение сборки}

\begin{itemize}
  \item Используйте BuildKit (по умолчанию в современных версиях).
  \item Кэшируйте пакетные менеджеры через
        \flag{--mount=type=cache}.
  \item Копируйте манифесты зависимостей раньше кода.
  \item Исключайте лишнее через \file{.dockerignore}.
  \item Разделяйте этапы: сборка зависимостей, сборка приложения,
        упаковка.
  \item В CI используйте удалённый кэш и параллельные этапы.
\end{itemize}

\begin{code}[language=bash]
$ docker buildx build --progress=plain -t api:1.0 .
$ time docker build -q -t api:1.0 .
\end{code}

\section{Что не стоит делать}

\begin{itemize}
  \item \textbf{Не используйте \flag{--squash} как основной приём.}
        Он ломает переиспользование слоёв и усложняет кэш.
  \item \textbf{Не удаляйте файлы в отдельных слоях} в надежде уменьшить
        образ.
  \item \textbf{Не гонитесь за минимальным размером ценой ремонтопригодности.}
        Distroless без shell усложняет отладку.
  \item \textbf{Не устанавливайте пакеты без фиксации версий}~---
        сборка перестанет быть воспроизводимой.
  \item \textbf{Не храните в образе секреты}~--- их не убрать удалением
        файла.
\end{itemize}

\begin{summarybox}
\begin{itemize}
  \item Место занимают базовый образ, кэши, инструменты сборки
        и лишние файлы.
  \item Минимальная база и многоэтапная сборка дают наибольший эффект.
  \item Чистка должна происходить в том же слое, где создан мусор.
  \item Кэш-монтирования и удалённый кэш ускоряют сборку, не увеличивая
        образ.
  \item Измеряйте результат \cmd{docker history}, \cmd{dive}
        и \cmd{docker image ls}.
\end{itemize}
\end{summarybox}

\section*{Вопросы для самопроверки}
\begin{selfcheck}
  \item Почему \cmd{rm} в отдельном слое не уменьшает образ?
  \item Чем distroless отличается от alpine и когда он уместен?
  \item Какие флаги Go уменьшают размер бинарника?
  \item Зачем нужен \flag{--mount=type=cache}?
  \item Почему \file{.dockerignore} влияет и на размер, и на безопасность?
  \item Что измерять, чтобы понять, какой слой самый тяжёлый?
\end{selfcheck}
